\section{可辨微观状态与任务相对宏观表示}
\label{sec:operational}

对有限集 $S,T$，以 $\Stoch(S,T)$ 记从 $S$ 到 $T$ 的行随机 Markov
核；通道按从左到右的次序复合。映射 $q:S\to T$ 所诱导的确定性核
记为 $Q_q$。全变差、复合和行向量约定集中列于附录~\ref{app:markov-tv}。

\subsection{有限观测--动力学呈现}

设 $X$ 为非空有限集合，$\cK\subseteq\Stoch(X,X)$ 为非空 Markov 核族。给定满射
\[
u:X\twoheadrightarrow U,
\]
称为\emph{微观探针}：它汇总当前系统边界内事先声明的微观观测。给定非恒定任务
\[
g=t\circ u:X\twoheadrightarrow O.
\]
四元组 $(X,\cK,u,g)$ 称为一个有限观测--动力学呈现。$u$ 和 $\cK$ 都属于问题设定；它们共同确定当前模型中的微观可辨关系。

从 $\Pi_0=\ker u$ 开始构造分区。已知 $\Pi_n$ 后，只在每个旧块内部继续区分状态：$x,x'$ 在 $\Pi_{n+1}$ 中同块，当且仅当它们在 $\Pi_n$ 中同块，且对每个 $K\in\cK$、每个 $C\in\Pi_n$，
\[
K(x,C)=K(x',C).
\]
分区至多严格细化 $|X|-1$ 次，故稳定于 $\Pi_\infty$。记
\[
p:X\twoheadrightarrow B:=X/\Pi_\infty.
\]
称 $B$ 为该呈现的\emph{可辨微观空间}。

\begin{proposition}[可辨微观基线]
\label{prop:operational-baseline}
上述构造满足：
\begin{enumerate}
  \item 存在唯一的 $\bar u:B\to U$ 与 $\bar g:B\to O$，使 $u=\bar u\circ p$、$g=\bar g\circ p$；
  \item 对每个 $K\in\cK$，存在唯一 $\bar K\in\Stoch(B,B)$，使
  \[
  KQ_p=Q_p\bar K;
  \]
  \item 在同时满足前两项的确定性商中，$p$ 最粗：若 $q:X\twoheadrightarrow Z$ 保留 $u$ 且使每个 $K$ 闭合，则存在唯一满射 $h:Z\twoheadrightarrow B$ 使 $p=h\circ q$。
\end{enumerate}
\end{proposition}

命题的证明是稳定分区细化的标准有限论证，见附录~\ref{app:operational-baseline}。空间 $B$ 相对于预先给定的 $(u,\cK)$ 定义；它保留探针当前可见或可由允许动力学在未来放大的差异，并合并其余精确冗余。

\subsection{先比较信息，再比较状态数}

从现在起在规范空间 $B$ 上讨论宏观状态。一个\emph{任务保真宏观表示}是满射
\[
s:B\twoheadrightarrow Y,
\qquad
\bar g=r\circ s
\]
其中 $r:Y\to O$。若 $s$ 非单射，则它是严格信息商，遗忘了至少一个可辨微观差异。等价地，在有限函数空间中有
\[
\bar g^{*}\mathbb R^O
\subseteq
s^{*}\mathbb R^Y
\subsetneq
\mathbb R^B.
\]
本文用 $|Y|/|B|$ 衡量这种已经由信息关系认证的压缩，而不直接使用可被编码填充任意放大的 $|X|$。

\begin{proposition}[严格复制填充不变性]
\label{prop:padding}
若把每个 $x\in X$ 替换成任意多个复制 $(x,a)$，并令探针、任务和全部动力学都忽略复制标签且使该标签与未来可观测分布无关，则填充后构造出的可辨微观空间与 $B$ 规范同构；诱导核、任务以及所有后续候选 $s$ 的闭合缺陷不变。
\end{proposition}

完整构造见附录~\ref{app:operational-baseline}。该命题处理\emph{精确}冗余填充；近似冗余变量需要另行给出统计分辨率或近似等价阈值。

从此，复杂度的分母固定为 $|B|$。下一节检验当前宏观状态是否足以给出自身的下一步分布，并证明可辨约化严格保持闭合误差。
